\subsection{Гармоническое движение}
%%%%%%%%%%%%%%%%%%%%%%%%%%%%%%%%%%%%%%%%%%%%%%%%%%%%
%%%%%%%%%%%%%%%%%%%%%%%%%%%%%%%%%%%%%%%%%%%%%%%%%%%%
 \z{Определите время полёта камня от одного полюса Земли до другого по прямому тоннелю, прорытому через её центр. Плотность Земли считать постоянной, радиус $R_\text{з} = 6400$ км.}
%
{$t = 42$ мин.}
%%%%%%%%%%%%%%%%%%%%%%%%%%%%%%%%%%%%%%%%%%%%%%%%%%%% 
 \z[{\begin{wrapfigure}[5]{r}{0.15\textwidth}%
		\vspace{-20pt}%
		\begin{flushright}%
 			\includegraphics[width=0.15\textwidth]{10.4.2}%
		\end{flushright}%
	\end{wrapfigure}%
	}]%
	{Тело массы $m$, подвешенное на пружине жёсткости $k$, лежит на подставке. Подставку мгновенно убирают. Опишите движение тела, если первоначально пружина:
	\pu не деформирована;
	\pu сжата и её деформация равна $l$.
	}
%
{\puan $x = \frac{mg}{k}(\cos \omega t - 1);$ ~ \puan $x = \left(\frac{mg}{k} + l\right)(\cos \omega t - 1)$.}
%%%%%%%%%%%%%%%%%%%%%%%%%%%%%%%%%%%%%%%%%%%%%%%%%%%% 
 \z[{\begin{wrapfigure}[5]{r}{0.3\textwidth}%
		\vspace{-10pt}%
		\begin{flushright}%
 			\includegraphics[width=0.3\textwidth]{10.4.3}%
		\end{flushright}%
	\end{wrapfigure}%
	}]%
	{По гладкой горизонтальной поверхности стола движутся с постоянной скоростью $v$ два бруска массами $m$ и $3m$, связанные нитью. Между брусками находится пружина жёсткостью $k$, сжатая на величину $x_0$ (см. рисунок). Пружина прикреплена только к бруску массой $m$. Размеры брусков малы по сравнению с длиной нити, массой пружины пренебречь, скорость брусков направлена вдоль нити. Во время движения нить обрывается, и бруски разъезжаются вдоль начального направления нити.
	\lpu Найти скорость бруска массой $3m$ после его отделения от пружины. \lpuend
	\lpu Найти время соприкосновения пружины с бруском массой $3m$, считая от момента разрыва нити. \lpuend
	}
%
{\puan $v_2 = v + \frac{x_0}{2}\sqrt{\frac{k}{3m}}$; ~ \puan $t = \frac{\pi}{4}\sqrt{\frac{3m}{k}}$}
%%%%%%%%%%%%%%%%%%%%%%%%%%%%%%%%%%%%%%%%%%%%%%%%%%%%	
 \z{По гладкой горизонтальной плоскости со скоростью $v$ скользит тонкий однородный брусок длины $l$. Брусок наезжает на обширный шероховатый участок плоскости. Через какое время брусок остановится, если коэффициент трения равен $\mu$?}
%
{$t = \frac{\pi}{2}\sqrt{\dfrac{l}{\mu g}} \text{при } v \leq \sqrt{\mu g l}$;~ $t =\frac{v'}{\mu g} + \sqrt{\frac{l}{\mu g}} \arccos \frac{v'}{v}$,     
    где $v' = \sqrt{v^2 - \mu g l}$.}
%%%%%%%%%%%%%%%%%%%%%%%%%%%%%%%%%%%%%%%%%%%%%%%%%%%%
\z[{\begin{wrapfigure}[5]{r}{0.2\textwidth}%
		\vspace{-10pt}%
		\begin{flushright}%
 			\includegraphics[width=0.2\textwidth]{10.4.5}%
		\end{flushright}%
	\end{wrapfigure}%
	}]%
	{Гладкую однородную верёвку длины $l$ удерживают в вертикальном колене изогнутой трубы так, что нижний конец её касается горизонтальной части трубы. Верёвку отпускают. Через какое время она полностью окажется в горизонтальном колене? Трением пренебречь. Как изменится это время, если вначале часть верёвки уже находилась в горизонтальном колене?}
%
{$t = \dfrac{\pi}{2}\sqrt{\frac{l}{g}}$, не изменится}